% Volume III, Paper 5 (drafted as "Chapter 31")
\chapter{Chamber Games and Higher-Order Interactions: Pairwise Reductions, Interfaces, and Entrainment}\label{chap:V3P5}
\chaptermark{Chamber Games}
\begin{chapterabstract}
The proposed programme places a three-player public-goods game on every chamber of an $\Atwo$
complex, with strategies updated by the Fermi rule. It asks for four claims. (a) A ``topological
shielding cage'' built from the $\PG(2,q)$ flags lowers the critical synergy to
$r_c=1+2/(q+1)$. (b) This is far below the value $r_c>3$ of ``traditional regular networks'',
a ``topological catalysis'' that grows with the thickness $q$. (c) A $\Z/3$-chiral
rock--paper--scissors variant has a torus attractor and persistent oscillations in the
replicator dynamics. (d) There is a global phase diagram in $|F|/|V|$. We prove the
following.

(1) \emph{Pairwise reduction.} The public-goods payoff is linear in the strategies. On
\emph{every} $2$-complex the chamber game is therefore the additive pairwise game
$\pi_i=(\frac r3-1)D_2(i)s_i+\frac r3\sum_jw_{ij}s_j$ on the $1$-skeleton, where $w_{ij}$ is
the number of chambers on the edge $ij$. Under (L), $w\equiv q+1$. After normalization by
$D_2$, the flag geometry disappears entirely, and only the degree $2\kappa$ of the
$1$-skeleton remains. There is no higher-order mechanism to exploit.

(2) \emph{Thresholds.} Defection is strictly dominant in every configuration for $r<3$. The
mean-field (replicator/Fermi) threshold is exactly $r_c=3$, for every $q$ and every $|F|/|V|$.
At zero noise, a cooperator cluster $A$ is stable iff $r>3\kappa/(\kappa+\delta(A))$, where
$\delta(A)$ is an integer boundary margin. So $5/3=1+2/(q+1)$ at $q=2$ is not the threshold of
any cluster on any $q=2$ complex. A single mutant needs $r>3\kappa/(\kappa-1)$ ($=3.5$ for
$q=2$). The best natural cage, the vertex star of $Y_{168}$, gives $7/3$. On $Y_{21}=K_{7,7,7}$
no cluster resists below $r=3$.

(3) \emph{No catalysis.} In Monte Carlo runs at noise $K=0.1$, $Y_{168}$ behaves like a random
$14$-regular graph: cooperators die at $r=2.7$, and the transition is near $3$. The
\emph{thin} complex, the triangular lattice $T_{12}$, keeps $40\%$ cooperators at $r=2.7$.
Thickness makes the complex an expander, and expanders are mean-field. The effect claimed in
(b) is reversed.

(4) \emph{Cyclic dominance.} For the three-strategy replicator dynamics,
$\frac{d}{dt}\log(x_1x_2x_3)=(a-b)(1-3\sigma)$. So there is global convergence ($a>b$), a
heteroclinic attractor ($a<b$), or a neutral family of closed orbits ($a=b$). There are never
limit cycles, and never a torus, since the phase space is a triangle. With mutation rate
$\mu$ a Hopf bifurcation occurs at $\mu_H=(b-a)/18$. Below it there is a stable limit cycle
(numerically supercritical): this is the correct form of ``persistent oscillations''.

(5) \emph{Nonlinear and chiral chamber games (Sections~\ref{V3P5:sec:nonlinear},
\ref{V3P5:sec:chiral}).} A chamber game is pairwise iff its synergy is affine. Symmetric
chamber games are exact potential games, so best-response learning selects a state that does not
depend on the geometry. Role-dependent (chiral) games break the potential. In the
supply-chain game with hold-up, logit learning in the mean field has a Hopf bifurcation at noise
$K_H=\frac14$ and periodic orbits below it. On large sparse complexes, however, the
oscillation of class densities decays with size. Under imitation, flat walls of the thin
lattice give the threshold $\frac{18}{11}$, while thick stars give $\frac{63}{19}$. We find no
spectral-gap law, and we show that the economic ``iff' proposed for dominant voluntary
cooperation is false.
\end{chapterabstract}

\section*{Status of the results}

Every numbered statement is proved, with two exceptions. The finite-noise thresholds are
Monte Carlo observations (Remark~\ref{V3P5:rem:mc}). In Theorem~\ref{V3P5:thm:mutation} the
uniqueness of the interior equilibrium and the supercriticality are verified numerically. All
checks are in \texttt{code/v3p5\_verify.py} (Table~\ref{V3P5:tab:battery}). Notation follows
Chapters~\ref{chap:V3P1}--\ref{chap:V3P4}: $Y$ satisfies (L), $\kappa=q^{2}+q+1$, the $1$-skeleton is $2\kappa$-regular,
every vertex lies in $D_2=\kappa(q+1)$ chambers, and $|F|/|V|=\kappa(q+1)/3$.

\section{The model and its pairwise reduction}\label{V3P5:sec:model}

\begin{definition}[Chamber public-goods game]\label{V3P5:def:pgg}
A configuration is $s\in\{0,1\}^{V}$, with $1=C$. In a chamber $\Delta$ each cooperator pays
$1$; the pool is multiplied by $r$ and shared equally among the three members. The payoff of
$i$ is
$\pi_i(s)=\sum_{\Delta\ni i}\bigl(\frac r3\,n_C(\Delta)-s_i\bigr)$, and the normalized payoff is
$\hat\pi_i=\pi_i/D_2(i)$. Under the \emph{Fermi rule} \cite{ST98} a uniformly random vertex $i$
picks a uniformly random neighbour $j$ and adopts $s_j$ with probability
$f(\hat\pi_i-\hat\pi_j)$, where $f(z)=1/(1+e^{z/K})$. The master equation is
\[
\frac{d}{dt}P(s)=\sum_{i}\bigl[w_i(s^{(i)})P(s^{(i)})-w_i(s)P(s)\bigr],\qquad
w_i(s)=\frac1{n\deg i}\sum_{j\sim i,\ s_j\ne s_i}f(\hat\pi_i-\hat\pi_j),
\]
where $s^{(i)}$ is $s$ with $s_i$ flipped.
\end{definition}

\begin{theorem}[Pairwise reduction]\label{V3P5:thm:reduction}
On any finite $2$-complex,
\[
\pi_i=\Bigl(\frac r3-1\Bigr)D_2(i)\,s_i+\frac r3\sum_{j}w_{ij}s_j,\qquad w_{ij}=\#\{\Delta\in F:\ i,j\in\Delta\}.
\]
So the chamber game is an additive pairwise game on the weighted $1$-skeleton. Under (L),
$w_{ij}=(q+1)[i\sim j]$, $D_2=\kappa(q+1)$, and
\[
\hat\pi_i=\Bigl(\frac r3-1\Bigr)s_i+\frac{r}{3\kappa}N_C(i),\qquad N_C(i)=\#\{j\sim i:s_j=1\}.
\]
The normalized dynamics depends on the complex only through its $1$-skeleton.
\end{theorem}

\begin{proof}
$n_C(\Delta)=\sum_{j\in\Delta}s_j$. Summing $\frac r3n_C(\Delta)$ over the chambers at $i$
counts $s_i$ $D_2(i)$ times and $s_j$ $w_{ij}$ times. Under (L), each edge lies in $q+1$
chambers.
\end{proof}

\begin{remark}\label{V3P5:rem:noflags}
The ``shielding cage from the $\PG(2,q)$ flags'' has no separate existence. The flags enter
only through the constant $q+1$, which cancels after normalization. Any two complexes with the
same $1$-skeleton and constant $w$ have identical dynamics, whether or not they satisfy the
flag condition. This is a structural feature of public goods, which are linear. A nonlinear
group payoff (threshold or volunteer games) would be genuinely higher-order.
\end{remark}

\begin{proposition}[Dominance and the mean-field threshold]\label{V3P5:prop:mf}
\begin{enumerate}
\item[(i)] In every configuration, switching $i$ from $D$ to $C$ changes $\pi_i$ by
$D_2(r/3-1)$. So $D$ is strictly dominant for $r<3$ and $C$ for $r>3$.
\item[(ii)] In the well-mixed (mean-field) limit the Fermi dynamics gives
$\dot x=x(1-x)\tanh\bigl(\frac{r/3-1}{2K}\bigr)$, and the replicator dynamics gives
$\dot x=x(1-x)(r/3-1)$. Both have the threshold $r_c=3$, independent of $q$ and $|F|/|V|$.
\end{enumerate}
\end{proposition}

\begin{proof}
(i) is read off Theorem~\ref{V3P5:thm:reduction}. (ii) With independent neighbours,
$\mathbb E(\hat\pi_C-\hat\pi_D)=r/3-1$. The rates $D\to C$ and $C\to D$ are
$x(1-x)f(-\Delta)$ and $x(1-x)f(\Delta)$, and $f(-\Delta)-f(\Delta)=\tanh(\Delta/2K)$.
\end{proof}

\section{Zero noise: pairs, clusters, cages}\label{V3P5:sec:zero}

As $K\to0$ the Fermi rule becomes ``imitate a strictly better neighbour''.

\begin{lemma}[Pair criterion]\label{V3P5:lem:pair}
For adjacent $i$ ($C$) and $j$ ($D$), $\hat\pi_i>\hat\pi_j$ iff
$\kappa+N_C(i)-N_C(j)>0$ and $r>3\kappa/(\kappa+N_C(i)-N_C(j))$. Moreover
$N_C(i)-N_C(j)\le2\kappa-2-c_{ij}\le2\kappa-q-3$, where $c_{ij}\ge q+1$ is the number of
common neighbours.
\end{lemma}

\begin{proof}
$\hat\pi_i-\hat\pi_j=\frac r3-1+\frac r{3\kappa}(N_C(i)-N_C(j))$. Let $C_{ij}$ be the set of
common neighbours and $m$ the number of cooperators in it. Since $s_j=0$, $N_C(i)\le
m+(2\kappa-1-c_{ij})$. Since $s_i=1$, $N_C(j)\ge m+1$. The edge $ij$ lies in $q+1$ chambers,
whose third vertices are common neighbours.
\end{proof}

For a set $A\subsetneq V$ of cooperators, define the \emph{boundary margin}
$\delta(A)=\min\{N_A(i)-N_A(j):i\in A,\ j\notin A,\ i\sim j\}$, where $N_A(x)$ is the number
of neighbours of $x$ in $A$.

\begin{theorem}[Cluster stability at zero noise]\label{V3P5:thm:cluster}
Under the zero-noise rule, no cooperator of $A$ ever imitates a defector iff
$\kappa+\delta(A)>0$ and $r>r_A:=3\kappa/(\kappa+\delta(A))$. In particular:
\begin{enumerate}
\item[(i)] a cluster can resist for some $r<3$ only if $\delta(A)>0$;
\item[(ii)] a single defector in a sea of cooperators, and a single cooperator in a sea of
defectors, both change sides at $r_*=3\kappa/(\kappa-1)=3+\frac3{q(q+1)}$. This is $3.5$ for
$q=2$, and it is \emph{above} the mean-field value;
\item[(iii)] every zero-noise cluster threshold on a $q=2$ complex has the form $21/(7+\delta)$
with $\delta\in\Z$. Since $21/(7+\delta)=5/3$ would need $\delta=5.6$, the value
$1+2/(q+1)=5/3$ is not the threshold of any cluster.
\end{enumerate}
\end{theorem}

\begin{proof}
Apply Lemma~\ref{V3P5:lem:pair} to every boundary pair. For (ii), use the sets $V\setminus\{v\}$
($\delta=-1$) and $\{v\}$: a lone $C$ beats a $D$ neighbour with one cooperating neighbour iff
$\frac r3-1>\frac r{3\kappa}$.
\end{proof}

\begin{proposition}[$Y_{21}$ has no cages]\label{V3P5:prop:Y21}
On $Y_{21}$ every proper nonempty $A$ has $\delta(A)\le0$, with equality iff $|A\cap V_t|$ is
the same for all three types. So $r_A\ge3$, and no cluster resists below the mean-field
threshold.
\end{proposition}

\begin{proof}
The $1$-skeleton is $K_{7,7,7}$ (Proposition~\ref{V3P2:prop:k777}), so
$N_A(x)=|A|-a_{\tau(x)}$ with $a_t=|A\cap V_t|$, and $\delta(A)=\min(a_{t'}-a_t)$ over
$t\ne t'$ with $a_t>0$, $a_{t'}<7$. Let $a_{t_0}=\max a_t\ge1$. If some $a_{t'}<a_{t_0}$, the
pair $(t_0,t')$ gives $\le-1$. Otherwise all $a_t$ are equal and $<7$, and every pair gives
$0$.
\end{proof}

\begin{remark}[Natural cages on the covers]\label{V3P5:rem:cages}
Computed margins: vertex stars have $\delta=-6,0,2$ on $Y_{21},Y_{84},Y_{168}$; chambers have
$\delta=0$ and edges $\delta=-1$ everywhere. The star of $Y_{168}$ is stable exactly for
$r>21/9=7/3$, which is checked directly. On the thin torus $T_{12}$ ($\kappa=3$), stars have
$\delta=1$ ($r>9/4$) and a half-torus strip has $\delta=2$ ($r>9/5$). These are statements
about zero noise only.
\end{remark}

\section{Finite noise: no topological catalysis}\label{V3P5:sec:mc}

\begin{remark}[Monte Carlo]\label{V3P5:rem:mc}
We ran the Fermi dynamics with normalized payoffs, starting from $50\%$ cooperators, with $60$
replicas and $400$ Monte Carlo steps. The table gives the final cooperator fraction.
\begin{center}
\small
\begin{tabular}{llcccccc}
\toprule
$K$ & structure & $r=2.4$ & $2.7$ & $2.9$ & $3.0$ & $3.1$ & $3.3$\\
\midrule
$0.5$ & $Y_{168}$ & 0.00 & 0.00 & 0.00 & 0.29 & 0.95 & 1.00\\
& $Y_{84}$ & 0.00 & 0.00 & 0.05 & 0.27 & 0.80 & 1.00\\
& random $14$-regular & 0.00 & 0.00 & 0.00 & 0.17 & 0.96 & 1.00\\
& $T_{12}$ (thin) & 0.00 & 0.00 & 0.09 & 0.64 & 0.93 & 1.00\\
\midrule
$0.1$ & $Y_{168}$ & 0.00 & 0.00 & 0.31 & 0.80 & 1.00 & 1.00\\
& $Y_{84}$ & 0.00 & 0.00 & 0.00 & 0.50 & 0.97 & 1.00\\
& random $14$-regular & 0.00 & 0.00 & 0.01 & 0.76 & 1.00 & 1.00\\
& $T_{12}$ (thin) & 0.00 & 0.40 & 0.69 & 0.79 & 0.89 & 1.00\\
\bottomrule
\end{tabular}
\end{center}
The rows for $K=0.1$ and $Y_{168}$, $T_{12}$ and the random graph are reproduced by the
battery. The $K=0.5$ rows and $Y_{84}$ come from the same code in an exploratory run. The
thick complexes are indistinguishable from a random $14$-regular graph, and their transition is
near the mean-field value $3$. The thin (Euclidean) complex supports cooperation at lower $r$:
this is the classical network reciprocity of lattices \cite{NM92,SPS09}. The ``topological
catalysis'' of the proposed programme is reversed. The $1$-skeleton of a thick complex is an
expander (Chapter~\ref{chap:V3P1}), its neighbourhoods are large ($2\kappa=14$), and the dynamics is close to
mean field. The sizes are small ($n\le168$), and we claim no critical values beyond this table.
\end{remark}

\begin{remark}[The comparison value ``$r_c>3$'']\label{V3P5:rem:three}
For groups of size $G$ the well-mixed threshold is $r=G$, here $3$. On lattices with
overlapping groups, network reciprocity gives $r_c<G$; for example $r_c/G\approx0.75$ on the
square lattice with $G=5$ \cite{SPS09}. ``$r_c>3$ on traditional regular networks'' is not the
correct baseline: our random $14$-regular graph has its transition near $3$.
\end{remark}

\paragraph{The ``global phase diagram''.} Under (L), $|F|/|V|=\kappa(q+1)/3$ is a function of
$q$ ($7$ for $q=2$). The proved thresholds are the mean-field $r_c=3$ for all $q$ and the
single-mutant $r_*=3+3/(q(q+1))$; both come from the table in line (E) of the battery. The
normalized dynamics does not contain $|F|/|V|$ at all (Theorem~\ref{V3P5:thm:reduction}). A
phase diagram in $|F|/|V|$ therefore carries no information beyond $q$, and no finite-noise
phase boundary is proved.

\section{Cyclic dominance and oscillations}\label{V3P5:sec:rps}

Give each vertex a strategy in $\Z/3$. In every chamber it plays the rock--paper--scissors
matrix $A=\begin{psmallmatrix}0&a&-b\\-b&0&a\\a&-b&0\end{psmallmatrix}$ against the two other
members: $k$ beats $k+1$ (gain $a$) and loses to $k-1$ (loss $b$). The chiral asymmetry is
$a\ne b$. By the argument of Theorem~\ref{V3P5:thm:reduction} the chamber game is again
$(q+1)$ times a pairwise game, and its mean field is the replicator equation
$\dot x_i=x_i((Ax)_i-x^{t}Ax)$ on the triangle $\Delta^{2}$.

\begin{theorem}[No limit cycles, no torus]\label{V3P5:thm:rps}
Let $V=x_1x_2x_3$ and $\sigma=x_1x_2+x_2x_3+x_3x_1\le\frac13$. Along the replicator flow,
$\frac{d}{dt}\log V=(a-b)(1-3\sigma)$. Hence:
\begin{enumerate}
\item[(i)] if $a>b$, every interior orbit converges to the barycentre;
\item[(ii)] if $a<b$, every interior orbit other than the barycentre converges to the
heteroclinic boundary cycle, and its oscillations slow down without bound;
\item[(iii)] if $a=b$, $V$ is conserved and interior orbits are closed. This family is neutral,
not an attractor.
\end{enumerate}
In no case is there an isolated periodic orbit. Since the phase space is two-dimensional, there
is no invariant torus.
\end{theorem}

\begin{proof}
$\frac{d}{dt}\log V=\sum_i(Ax)_i-3x^{t}Ax$. The column sums of $A$ are $a-b$, and
$x^{t}Ax=(a-b)\sigma$. Also $\sigma\le(\sum x_i)^{2}/3$, with equality only at the barycentre.
In (i) and (ii), $V$ is a strict Lyapunov function on the interior minus the barycentre; LaSalle
\cite{LaSalle60} gives the limits. In (iii) the level sets of $V$ are closed curves. A torus
does not embed in the plane.
\end{proof}

\begin{theorem}[Mutation creates a limit cycle]\label{V3P5:thm:mutation}
Consider $\dot x_i=x_i((Ax)_i-x^{t}Ax)+\mu(1-3x_i)$ with $\mu>0$. At the barycentre the
linearization has eigenvalues $\frac{b-a}6-3\mu\pm i\frac{\sqrt3(a+b)}6$. For $a<b$ a Hopf
bifurcation occurs at $\mu_H=(b-a)/18$. For $0<\mu<\mu_H$, the triangle is forward invariant and
the barycentre is a repelling focus. So, if the barycentre is the only equilibrium, every other
orbit accumulates on a periodic orbit by Poincaré--Bendixson.
\end{theorem}

\begin{proof}
On $\{\sum\xi=0\}$ the linearization is $\frac13A-3\mu I$, because $\delta(x^{t}Ax)=0$ there.
The eigenvalues of the circulant $A$ on that plane are $a\omega-b\omega^{2}$ and its conjugate,
i.e.\ $\frac{b-a}2\pm i\frac{\sqrt3(a+b)}2$. On the face $x_i=0$ we have $\dot x_i=\mu>0$.
\end{proof}

Numerically, for $(a,b)=(1,2)$ the barycentre is the only equilibrium, on a $400\times400$ grid
and for four parameter sets. The cycle is reached from inside and from near the boundary, with
amplitudes $0.397,0.160,0.048$ at $\mu/\mu_H=0.5,0.9,0.99$. The squared amplitude scales like
$1-\mu/\mu_H$, which indicates a supercritical bifurcation, and the cycle is gone at
$1.05\mu_H$. This agrees with \cite{Mob10}. It is a limit cycle, a one-dimensional attractor,
not a torus. A quasi-periodic torus would need at least three dimensions, for example the three
type classes as separate populations; we do not treat that case.

\section{Nonlinear chamber games}\label{V3P5:sec:nonlinear}

This section treats open problem (d) of Remark~\ref{V3P5:rem:scope}. Let $f\colon\{0,1,2,3\}\to\R$
with $f(0)=0$ be the synergy of a chamber. Every chamber pays each member $f(n_C)/3$ and charges
each cooperator $1$, so
\[
\pi_i(s)=\sum_{\Delta\ni i}\Bigl(\tfrac13f(n_C(\Delta))-s_i\Bigr).
\]
The examples are the linear game $f=r(0,1,2,3)$, the threshold game $f=r(0,0,2,3)$ (so
$f(0)=f(1)=0$, $f(2)=2r$, $f(3)=3r$), and the quadratic game $f(n)=rn^{2}$. On $\{0,1\}^{3}$,
\[
f(s_1+s_2+s_3)=c_1(s_1+s_2+s_3)+c_2(s_1s_2+s_1s_3+s_2s_3)+c_3s_1s_2s_3,
\]
with $c_1=f(1)$, $c_2=f(2)-2f(1)$ and $c_3=f(3)-3f(2)+3f(1)$. The threshold game has
$(c_1,c_2,c_3)=(0,2r,-3r)$.

\begin{theorem}[Irreducibility]\label{V3P5:thm:irreducible}
Call a game \emph{pairwise} if every $\pi_i$ is a sum of a function of $s_i$ and functions of
$(s_i,s_j)$ over the neighbours $j$ of $i$. On a $2$-complex in which every vertex lies in a
chamber, the chamber game is pairwise iff $c_2=c_3=0$, i.e.\ iff $f$ is affine on
$\{0,1,2,3\}$. Otherwise $\pi_i$ contains the monomials $\frac{c_2}3s_js_k$ and
$\frac{c_3}3s_is_js_k$ for every chamber $\{i,j,k\}$, which are genuine three-body terms.
\end{theorem}

\begin{proof}
Multilinear expansions of functions on $\{0,1\}^{V}$ are unique. A pairwise payoff contains only
the monomials $1$, $s_i$, $s_j$ and $s_is_j$ with $j\sim i$. The monomial $s_js_k$, for
$\{i,j,k\}\in F$, arises only from that chamber, with coefficient $c_2/3$. Likewise
$s_is_js_k$ arises only from that chamber, with coefficient $c_3/3$. For affine $f$ we are back in
Theorem~\ref{V3P5:thm:reduction}.
\end{proof}

\begin{theorem}[Exact potential]\label{V3P5:thm:potential}
For every $f$ and every $2$-complex, $\Phi(s)=\sum_{\Delta}\bigl(\frac13f(n_C(\Delta))-n_C(\Delta)\bigr)$
is an exact potential: switching $i$ changes $\pi_i$ and $\Phi$ by the same amount. Hence:
\begin{enumerate}
\item[(i)] the Nash equilibria are the local maxima of $\Phi$ under single switches;
\item[(ii)] logit (best-response) dynamics is reversible, with Gibbs measure $\propto e^{\Phi/K}$;
\item[(iii)] if $\psi(n)=\frac13f(n)-n$ is maximized over $\{0,1,2,3\}$ at $n\in\{0,3\}$ (true for
convex $f$ and for the threshold game), then as $K\to0$ the stochastically stable state is all-$D$
if $f(3)<9$ and all-$C$ if $f(3)>9$, \emph{on every complex}.
\end{enumerate}
In particular the threshold game switches at $r=3$, and $f=rn^{\gamma}$ switches at $r=3^{2-\gamma}$,
whatever the geometry.
\end{theorem}

\begin{proof}
The change of $\pi_i$ is $\sum_{\Delta\ni i}\bigl(\frac13(f(k_\Delta+1)-f(k_\Delta))-1\bigr)$, where
$k_\Delta$ counts the other cooperators, and this is also the change of $\Phi$. (ii) is standard
for exact potential games. (iii) Every chamber contributes at most $\max\psi$, and the uniform
states attain that bound in every chamber simultaneously. For the threshold game,
$\psi=(0,-1,\frac{2r}3-2,r-3)$, and $\psi(2)\le\max(0,r-3)$.
\end{proof}

\begin{proposition}[Type-stratified equilibria]\label{V3P5:prop:typeNash}
On any complex in which every chamber has one vertex of each type ($\Atwo$ complexes, thick or
thin), the state ``$D$ exactly on one type class'' puts two cooperators in every chamber. It is a
Nash equilibrium iff $f(2)-f(1)\ge3\ge f(3)-f(2)$. For the threshold game this is
$\frac32\le r\le3$. For $r\in(\frac32,3)$ it is a strict local maximum of $\Phi$, but not the
global one, so it is metastable under logit dynamics. Under imitation it is unstable: a
cooperator and a neighbouring defector see identical chambers, and the defector earns $D_2$
more.
\end{proposition}

\begin{proof}
A cooperator's other chamber members are one $C$ and one $D$, so switching changes its payoff by
$D_2(1-\frac13(f(2)-f(1)))$. A defector's other members are two $C$'s, so switching changes its
payoff by $D_2(\frac13(f(3)-f(2))-1)$.
\end{proof}

This is the one place where the $2$-cells matter for equilibria. What it uses is
partiteness (the $\Z/3$ grading), not $q$ and not the flag structure.

\paragraph{Imitation: mean field.} For the Fermi rule in the well-mixed limit,
$\dot x=x(1-x)\tanh(\Delta(x)/2K)$ with
$\Delta(x)=\frac13\sum_{k=0}^{2}\binom2kx^{k}(1-x)^{2-k}\bigl(f(k+1)-f(k)\bigr)-1$. For the threshold
game, $\Delta(x)=\frac r3(4x-3x^{2})-1$. All-$D$ is always stable. For $r>\frac94$ a stable mixed
state appears, and all-$C$ is stable iff $r>3$. So nonlinearity lowers the mean-field threshold
from $3$ to $\frac94$. This effect does not depend on the geometry.

\begin{proposition}[Star cages under imitation, $q\in\{1,2\}$]\label{V3P5:prop:starNL}
Let $A$ be the star of a vertex in a complex whose radius-$3$ balls are those of an $\Atwo$
building of order $q\in\{1,2\}$. Put $\varphi=f/r$. Then $A$ loses no cooperator at zero noise iff
$r>r_{\mathrm{star}}$, where
\[
r_{\mathrm{star}}(q)=\frac{3(q+1)(q^{2}+q+1)}{(q+1)\bigl(\varphi(3)+q\varphi(2)+q^{2}\varphi(1)\bigr)-\varphi(2)-2q\varphi(1)}.
\]
This gives $\frac73,\frac{63}{19},\frac{63}{55}$ for $q=2$ (linear, threshold, quadratic) and
$\frac94,\frac94,\frac9{11}$ for $q=1$.
\end{proposition}

\begin{proof}
A boundary vertex $u$ of the star has $q+1$ chambers through $v$ (three cooperators), $q(q+1)$
chambers with exactly one other cooperator, and $q^{2}(q+1)$ chambers with none. This follows from
the link: its neighbours in the star are the $q+1$ link-neighbours of $u$ in $\mathrm{Lk}(v)$, and
these are pairwise non-adjacent. The most dangerous outside neighbour $j$ is the third vertex of a
chamber over a boundary edge. It has one chamber with two cooperators and $2q$ chambers with one.
Then apply the zero-noise pair criterion. For $q=2$ the count at $j$, and the maximality of that
$j$, are verified exactly in the radius-$3$ ball of the building (a level-$12$ cover of the
canonical tower of Chapter~\ref{chap:V3P6}, whose ball sizes $1,14,98,560$ agree with those of the building) and
on $Y_{168}$. For $q=1$ they are verified on $T_{12}$ (line (NL4)).
\end{proof}

For large $q$ the displayed formula grows like $\frac32q$ for the threshold game and tends to $3$
in the other two cases, assuming the same counts. A linear $3$-uniform hypertree with the same
chamber degree $D_2$ gives $3D_2/(\varphi(3)+(D_2-2)\varphi(1))$ for the star, i.e.\ $21$ for the
threshold game when $D_2=21$. So the building's star beats the hypertree's star. The mechanism
is the $(q+1)$-fold overlap of chambers along edges, not the incidence ``two points on one line''.
The star does not beat the thin lattice, and it gets worse as $q$ grows.

\begin{remark}[Exhaustive comparison on $21$ vertices]\label{V3P5:rem:exhaustive}
For hypergraphs on $21$ vertices with $D_2=21$ we enumerated all $2^{21}-2$ cooperator sets and
computed the best zero-noise cage threshold $\min_A r_A$ (line (NL3)).
\begin{center}
\small
\begin{tabular}{lccc}
\toprule
hypergraph & linear & threshold & quadratic\\
\midrule
$Y_{21}$ (chambers) & $3$ & $\frac94$ & $0.677$\\
random tripartite, $3$ seeds & $2.739$ & $2.10$--$2.17$ & $0.670$--$0.685$\\
random $3$-uniform, $3$ seeds & $2.86$--$3$ & $2.17$--$2.25$ & $0.663$--$0.685$\\
\bottomrule
\end{tabular}
\end{center}
$Y_{21}$ is never better than the random models; for the linear game it is worse. The battery
reproduces seed $0$ of each random model. The quadratic cages are almost full cooperation. They
sit near the mean-field stability value $\frac35$ of all-$C$, where $\Delta(1)=\frac53r-1$.
\end{remark}

\begin{remark}[Monte Carlo]\label{V3P5:rem:nlmc}
Threshold game, Fermi rule, $K=0.1$, $50\%$ initial cooperators, $60$ replicas, $300$ sweeps.
Final cooperator fraction:
\begin{center}
\small
\begin{tabular}{lcccccc}
\toprule
structure & $r=2.25$ & $2.5$ & $2.75$ & $3.0$ & $3.25$ & $3.5$\\
\midrule
$Y_{21}$ & 0.00 & 0.00 & 0.50 & 0.77 & 0.85 & 0.90\\
$Y_{168}$ & 0.00 & 0.59 & 0.94 & 1.00 & 1.00 & 1.00\\
random, $n=168$, $D_2=21$ & 0.00 & 0.60 & 0.94 & 1.00 & 1.00 & 1.00\\
$T_{12}$ (thin) & 0.09 & 0.87 & 0.99 & 1.00 & 1.00 & 1.00\\
random, $n=144$, $D_2=6$ & 0.00 & 0.48 & 0.90 & 0.98 & 1.00 & 1.00\\
\bottomrule
\end{tabular}
\end{center}
For comparison, the linear game on the same structures (exploratory run) switches between $3.0$
and $3.25$. The exceptions are $T_{12}$, which already has $0.50$ at $2.75$, and $Y_{21}$, which
still has only $0.27$ at $3.25$. The battery reproduces the columns $2.25$--$3.0$ for $Y_{21}$, $T_{12}$,
$Y_{168}$ and the $D_2=21$ random hypergraph. Three facts stand out. Nonlinearity lowers the
transition from about $3$ to about $2.4$--$2.5$, in line with the mean-field value $\frac94$. The
thick complex $Y_{168}$ is indistinguishable from a random hypergraph with the same chamber
degree. The thin lattice is the most favourable structure, and the dense $Y_{21}=K_{7,7,7}$ the
least.
\end{remark}

\begin{remark}[Answer to the deepening request]\label{V3P5:rem:deep}
(1) ``Pairwise iff linear'' is true (Theorem~\ref{V3P5:thm:irreducible}); nonlinear chamber games
are genuinely three-body. (2) The incidence fact that is actually used is ``each edge lies in $q+1$
chambers''. Two edges at a vertex lie in at most one common chamber, not ``exactly one''. (3)
Under best-response learning the geometry is irrelevant for the selected state
(Theorem~\ref{V3P5:thm:potential}). The one structural effect, the type-stratified equilibria of
Proposition~\ref{V3P5:prop:typeNash}, comes from partiteness and is metastable. (4) Under
imitation, the $\Atwo$ star cage is much better than a hypertree star, but its threshold grows with
$q$ and exceeds the thin-lattice value (Proposition~\ref{V3P5:prop:starNL}). Exhaustively, $Y_{21}$
is no better than random hypergraphs with the same parameters (Remark~\ref{V3P5:rem:exhaustive}).
At finite noise the thick complex matches a random hypergraph (Remark~\ref{V3P5:rem:nlmc}). A
``topological shielding threshold significantly below that of random regular hypergraphs'' is
therefore not established, and the data contradict it. What nonlinearity does lower, uniformly,
is the mean-field threshold, from $3$ to $\frac94$.
\end{remark}

\section{Chiral games, interfaces, and an economic reading}\label{V3P5:sec:chiral}

Every chamber has one vertex of each type. Let the vertex of type $\rho$ play \emph{role} $\rho$ in
all its chambers, so that roles are ordered cyclically, $0\to1\to2\to0$. A \emph{chiral chamber
game} gives role $\rho$ the per-chamber payoff $u_\rho(s_0,s_1,s_2)=s_\rho\,g_\rho(s_{\rho-1},s_{\rho+1})$.
The gain $g_\rho$ may depend on the upstream member $\rho-1$ and the downstream member $\rho+1$
differently. A vertex's payoff is the average of $u_{\tau(i)}$ over its $D_2$ chambers. The
\emph{supply-chain game} has roles technology (0), manufacturing (1) and channel (2), with
\[
g_0=1-2s_2,\qquad g_1=2s_0-1,\qquad g_2=2s_1-1 .
\]
Manufacturing's investment pays only if technology invested; the channel's pays only if
manufacturing did; technology is held up once the channel invests. The \emph{three-body variant}
replaces $g_2$ by $4s_0s_1-1$: the channel's return is activated only when both upstream partners
invest.

\begin{theorem}[Potential criterion]\label{V3P5:thm:potcrit}
Let $Y$ satisfy \textup{(L)}, and for $\rho\ne\sigma$ let $D_{\rho\sigma}u_\rho$ be the mixed
difference of $u_\rho$ in $(s_\rho,s_\sigma)$, a function of the third strategy $s_\kappa$. The
network game has an exact potential iff $D_{\rho\sigma}u_\rho\equiv D_{\rho\sigma}u_\sigma$ for
all $\rho\ne\sigma$. It has a weighted potential iff there are $w_\rho>0$ with
$w_\rho D_{\rho\sigma}u_\rho\equiv w_\sigma D_{\rho\sigma}u_\sigma$. The symmetric games of
Section~\ref{V3P5:sec:nonlinear} satisfy the exact criterion. The supply-chain game and its
three-body variant satisfy neither. Their Monderer--Shapley circulation on an upstream edge is
$\pm(q+1)\cdot2/D_2$ ($\pm2/7$ on $Y_{21}$), and a single chamber has a $6$-cycle of strict
improvements
$(0,0,0)\to(1,0,0)\to(1,1,0)\to(1,1,1)\to(0,1,1)\to(0,0,1)\to(0,0,0)$. So there is not even an
ordinal potential.
\end{theorem}

\begin{proof}
A finite game has an exact potential iff every $2\times2$ subgame has zero circulation, i.e.\
$D_{ij}\pi_i=D_{ij}\pi_j$ for all players $i,j$ and all profiles of the others \cite{MS96}.
Players interact only along edges. For an edge $ij$ with roles $(\rho,\sigma)$, the $q+1$
chambers containing it all assign the same roles, and their third vertices are distinct. So
$D_{ij}\pi_i=\frac1{D_2}\sum_{\Delta\supset ij}D_{\rho\sigma}u_\rho(s_{k_\Delta})$. Setting all
third vertices to $0$, or all to $1$, gives necessity of the per-chamber identity; sufficiency is
clear. The weighted case is the same computation with $w_i\pi_i$. For the supply-chain game,
$D_{10}u_1=2$ while $D_{10}u_0=0$, which cannot be balanced by any positive weights. The
improvement cycle is read off the signs of $g_\rho$.
\end{proof}

In the well-mixed limit each type class is a population with cooperator fraction $x_\rho$, and
logit learning at noise $K$ gives
\begin{equation}\label{V3P5:eq:logitmf}
\dot x_\rho=\sigma\bigl(G_\rho(x)/K\bigr)-x_\rho,\qquad \sigma(z)=\frac1{1+e^{-z}},\qquad
G_\rho(x)=\mathbb E\,g_\rho .
\end{equation}

\begin{proposition}[Potential mean fields converge]\label{V3P5:prop:lyap}
If $w_\rho G_\rho=\partial\Phi/\partial x_\rho$ for some $\Phi$ and weights $w_\rho>0$, then
$L=\Phi/K+\sum_\rho w_\rho H(x_\rho)$, with $H(x)=-x\log x-(1-x)\log(1-x)$, increases strictly
along \eqref{V3P5:eq:logitmf} away from equilibria. So there are no periodic orbits.
\end{proposition}

\begin{proof}
$\dot L=\sum_\rho w_\rho\bigl(G_\rho/K-\operatorname{logit}x_\rho\bigr)\bigl(\sigma(G_\rho/K)-x_\rho\bigr)\ge0$,
because $\sigma$ is increasing. Equality holds iff every factor vanishes.
\end{proof}

\begin{theorem}[Endogenous cycles in cyclic chiral games]\label{V3P5:thm:cycle}
Suppose $G_\rho=G_\rho(x_{\rho-1})$ is strictly monotone with sign $\varepsilon_\rho$. Then
\eqref{V3P5:eq:logitmf} is a monotone cyclic feedback system on $[0,1]^{3}$.
\begin{enumerate}
\item[(i)] If $\varepsilon_0\varepsilon_1\varepsilon_2=+1$, the system is cooperative and
irreducible, and no nontrivial periodic orbit is asymptotically stable \cite{Smith95}.
\item[(ii)] If $\varepsilon_0\varepsilon_1\varepsilon_2=-1$, the equilibrium $x^{*}$ is unique.
Its linearization has eigenvalues $-1+|P|^{1/3}e^{i\pi(2k+1)/3}$, $k=0,1,2$, where
$P=\prod_\rho\sigma'\bigl(G_\rho(x^{*}_{\rho-1})/K\bigr)G_\rho'(x^{*}_{\rho-1})/K<0$. So $x^{*}$
is stable iff $|P|<8$. At $|P|=8$ a pair of eigenvalues crosses the imaginary axis at $\pm
i\sqrt3$. For $|P|>8$ a nonconstant periodic orbit exists (Mallet-Paret--Smith
\cite{MPS90}).
\end{enumerate}
For the supply-chain game $x^{*}=(\frac12,\frac12,\frac12)$ and $|P|=(2K)^{-3}$. The Hopf point is
$K_H=\frac14$, and the period there is $2\pi/\sqrt3=3.628$.
\end{theorem}

\begin{proof}
(i) With positive feedback the off-diagonal Jacobian entries are $\ge0$, and they form an
irreducible cyclic pattern in the interior. (ii) An equilibrium solves $x_\rho=F_\rho(x_{\rho-1})$
with $F_\rho=\sigma\circ(G_\rho/K)$ monotone of sign $\varepsilon_\rho$. The composite
$F_0\circ F_2\circ F_1$ is decreasing, so it has exactly one fixed point. The Jacobian is $-I+C$,
where $C$ is the cyclic matrix with entries $\sigma'G_\rho'/K$, and $C$ has eigenvalues $P^{1/3}$.
The real part of $-1+|P|^{1/3}e^{\pm i\pi/3}$ is $-1+|P|^{1/3}/2$. For $|P|>8$ the equilibrium is
unique and has a two-dimensional unstable manifold. The Poincar\'e--Bendixson theory for monotone
cyclic feedback systems \cite{MPS90} (quoted) then gives a nonconstant periodic orbit.
\end{proof}

\begin{remark}[Mean-field numerics]\label{V3P5:rem:cyclenum}
For the supply-chain game the cycle amplitude of $x_0$ is $0.294,0.287,0.263,0.206,0.098$ at
$K=0.05,0.10,0.15,0.20,0.24$, and $0$ for $K\ge0.26$. The period is $3.00,3.26,3.47,3.59,3.63$,
approaching $2\pi/\sqrt3$. The squared amplitude scales like $K_H-K$, which indicates a
supercritical bifurcation. The three-body variant is not a cyclic feedback system, since $x_2$
responds to both $x_0$ and $x_1$. Its Jacobian at $(\frac12,\frac12,\frac12)$ is
$-I+\frac1{4K}M$, where $M$ has characteristic polynomial $\lambda^{3}+4\lambda+8$. So its Hopf point
is $K_H=0.170582$, a quarter of the real part of the complex root. Numerically the
cycle exists below $K_H$ (amplitude $0.186$ at $K=0.05$) and not above. Its existence beyond the
local Hopf regime is not proved. Line (NC4) checks both Hopf points and the amplitudes.
\end{remark}

\begin{remark}[Networks: the cycle is a mean-field phenomenon]\label{V3P5:rem:cyclenet}
Logit dynamics of the supply-chain game on the complexes themselves, at $K=0.1$ (below $K_H$), gives
the following standard deviation of the class density $x_0(t)$ (exploratory runs over $120$ time
units):
\begin{center}
\small
\begin{tabular}{lcccccc}
\toprule
& $Y_{24}$ & $Y_{768}$ & $Y_{6144}$ & $T_{12}$ & $T_{24}$ & $T_{48}$\\
$n$ & $24$ & $768$ & $6144$ & $144$ & $576$ & $2304$\\
\midrule
std $x_0$, $K=0.1$ & $0.262$ & $0.091$ & $0.023$ & $0.094$ & $0.044$ & $0.025$\\
std $x_0$, $K=0.4$ & $0.224$ & $0.041$ & $0.015$ & $0.074$ & $0.036$ & $0.021$\\
\bottomrule
\end{tabular}
\end{center}
Here $Y_{768}$ and $Y_{6144}$ are levels $4$ and $5$ of the canonical $2$-tower of Chapter~\ref{chap:V3P6}. In the
mean field the standard deviation would be about $0.2$. On both thick and thin complexes it
decays with $n$. The single-site autocorrelation has only a weak negative lobe ($-0.06$ on
$Y_{768}$, $-0.09$ on $T_{24}$ at lag $3$), so local oscillations are weak and incoherent. The
limit cycle of Theorem~\ref{V3P5:thm:cycle} describes well-mixed populations; on $Y_{21}=K_{7,7,7}$
the drift is exactly mean-field. No coherent network rhythm was found on large sparse
complexes, thick or thin. Line (NC6) reproduces the decay for $Y_{24}\to Y_{168}$ and
$T_{12}\to T_{48}$.
\end{remark}

\begin{proposition}[Flat interfaces]\label{V3P5:prop:flat}
In the thin complex $T_m$ ($3\mid m$), let $A$ be a half-torus bounded by two walls parallel to a
lattice direction. At zero noise under imitation, $A$ loses no cooperator iff
\[
r>\frac{18}{3\varphi(3)+\varphi(2)-\varphi(1)},
\]
i.e.\ $r>\frac95$, $\frac{18}{11}$, $\frac35$ for the linear, threshold and quadratic games.
\end{proposition}

\begin{proof}
A boundary cooperator has $6$ chambers: three with three cooperators, two with two, and one with
one. Its defecting neighbour across the wall has one chamber with two cooperators, two with one,
and three with none. So the gap in normalized payoffs is
$r(3\varphi(3)+\varphi(2)-\varphi(1))/18-1$, with $D_2=6$. The two boundary vertices of each type are equivalent by translation.
\end{proof}

The flat wall is the exact mechanism behind the advantage of the thin lattice
(Remark~\ref{V3P5:rem:nlmc}). A boundary cooperator on a wall has only one ``lonely'' chamber out of
$6$. At the boundary of a thick star it has $q^{2}(q+1)=12$ out of $21$
(Proposition~\ref{V3P5:prop:starNL}). Enlarging the cluster to balls of radius $2$ and $3$ in the
thick building gives exactly the star thresholds again, as we computed in the level-$14$ cover. Thick
buildings have no flat interfaces: their spheres grow exponentially.

\begin{remark}[Why no analytic law $r_c(q+1,\lambda_2)$]\label{V3P5:rem:nolaw}
A finite-noise critical value is not sharply defined on finite complexes, where the dynamics is
absorbing. On the infinite building, the existence of a phase transition for this non-reversible
particle system is itself open. The data also rule out $\lambda_2$ as the controlling variable.
With normalized gaps $\lambda_2/\deg=0.089$ ($T_{12}$), $0.441$ (random, $D_2=6$), $0.601$
($Y_{168}$), $0.697$ (random, $D_2=21$) and $1.000$ ($Y_{21}$), the threshold-game transitions of
Remark~\ref{V3P5:rem:nlmc} lie at about $2.38$, $2.51$, $2.48$, $2.48$ and $2.75$. The ordering is
roughly monotone, but $Y_{168}$ sits below a hypergraph with a smaller gap. Interfaces are what the
proved statements identify (Propositions~\ref{V3P5:prop:starNL}, \ref{V3P5:prop:flat}). A
chamber-cluster approximation could be written down, but it would be an uncontrolled closure,
and we do not present its output as a law.
\end{remark}

\paragraph{An economic reading.} Firms (vertices) join three-party projects (chambers). Each
project has one firm of each role, and each firm has $D_2$ projects. Axioms:
\begin{itemize}
\item[(A1)] a firm chooses one policy, contribute ($1$) or not ($0$), for all its projects;
\item[(A2)] contributing costs $1$ per project;
\item[(A3)] a project's output depends on its contributors through a technology, which is
symmetric, $f(n_C)$ shared equally, or role-dependent, $u_\rho$;
\item[(A4)] payoffs are averaged over projects.
\end{itemize}
An R\&D consortium is the symmetric case, with a threshold technology when a breakthrough needs at
least two partners. A syndicated credit line is the same with ``lend'' for ``contribute''. A supply
chain with hold-up is the chiral case.

\begin{proposition}[Free riding, equilibrium and selection]\label{V3P5:prop:econ}
In the symmetric case:
\begin{enumerate}
\item[(i)] contributing is a dominant strategy iff $\min_{k=0,1,2}\bigl(f(k+1)-f(k)\bigr)\ge3$. In
particular it is \emph{never} dominant under a threshold technology ($f(1)=f(0)$);
\item[(ii)] universal contribution is a Nash equilibrium iff $f(3)-f(2)\ge3$;
\item[(iii)] under logit learning, universal contribution is selected as $K\to0$ iff $f(3)>9$
(when $\psi$ peaks at the ends), \emph{whatever the contract network}.
\end{enumerate}
The network matters only for non-equilibrium social learning (imitation), through its
interfaces (Proposition~\ref{V3P5:prop:flat}).
\end{proposition}

\begin{proof}
(i) A firm's gain from contributing in a project with $k$ other contributors is
$\frac13(f(k+1)-f(k))-1$; dominance must hold for every $k$ in every project. (ii) is the case
$k=2$. (iii) is Theorem~\ref{V3P5:thm:potential}(iii).
\end{proof}

The proposed programme asks to prove that voluntary cooperation is a globally dominant equilibrium
\emph{iff} the technology has a threshold and the network is moderately local. This is false in
both directions. A threshold technology rules out dominance (i). Equilibrium selection
under best-response learning ignores the network (iii). What locality does is help
\emph{imitation} spread cooperation across flat interfaces. In the chiral (supply-chain) case,
hold-up creates negative feedback and endogenous investment cycles for low noise. These are
coherent in a well-mixed market (Theorem~\ref{V3P5:thm:cycle}) but not across a large sparse
contract network (Remark~\ref{V3P5:rem:cyclenet}).

\section{Clocks, entrainment and anchors}\label{V3P5:sec:clocks}

The two obstacles found above are the decay of network oscillations
(Remark~\ref{V3P5:rem:cyclenet}) and the absence of flat interfaces in thick buildings
(Proposition~\ref{V3P5:prop:flat}). This section tests the two remedies proposed in the
programme. The first is a settlement clock carried by the gauge field of Chapter~\ref{chap:V3P3}. The second is
a committed core (``chain leader''). The first produces coherence, but by entrainment rather than
endogenously. The second cannot overcome the geometry.

\paragraph{The clock.} Each vertex carries a phase $\theta_i$ evolving by the flux sector of
Chapter~\ref{chap:V3P3}: identical frequencies $\omega$, coupling $K_1>0$, flux $2\pi\alpha$. Its settlement cycle
enters the supply-chain gain as $g_{\tau(i)}+\varepsilon\cos\theta_i$. The coupling is one-way:
strategies do not act on phases.

\begin{proposition}[The clock locks with a flux-selected twist]\label{V3P5:prop:clock}
The phases converge to an equilibrium of the flux sector in the rotating frame
(Theorem~\ref{V3P3:thm:gradient}). The ground states are the twisted states
$\theta_i=\omega t+c+\frac{2\pi m^{*}}3\tau(i)$ with $m^{*}\equiv-\operatorname{round}(\alpha)$
(Theorem~\ref{V3P3:thm:ground}). So a synchronized clock forces the three roles with a relative
phase lag of $\frac{2\pi m^{*}}3$, and the flux selects this chirality. From random initial phases
the clock reached the ground state, with order $>0.999$, on $Y_{24}$ and $Y_{168}$ for
$\alpha=0.3,1.2,-0.8$ (so $m^{*}=0,2,1$), and on $Y_{768}$. On the thin torus $T_{48}$ it stopped at
order $0.41$, trapped in a non-ground equilibrium. Expansion, which hurts cooperation, helps the
clock.
\end{proposition}

\begin{theorem}[Entrained cycle: linear response]\label{V3P5:thm:entrain}
Let $K>K_H=\frac14$, so that the mean-field equilibrium $x^{*}=(\frac12,\frac12,\frac12)$ of the
supply-chain game is a stable focus (Theorem~\ref{V3P5:thm:cycle}). Force it with a clock of
twist $m$, i.e.\ $\dot x_\rho=\sigma\bigl((G_\rho(x)+\varepsilon\cos(\omega t+\frac{2\pi
m\rho}3))/K\bigr)-x_\rho$. For small $\varepsilon$ this has a unique attracting periodic solution
$x^{*}+\varepsilon\operatorname{Re}\bigl(e^{i\omega t}X_m(\omega)\bigr)+O(\varepsilon^{2})$, with
\[
X_m(\omega)=\frac1{4K}\bigl(i\omega I+I-C\bigr)^{-1}f_m,\qquad f_m=\bigl(e^{2\pi im\rho/3}\bigr)_\rho,
\]
where $C$ is the cyclic matrix with entries $C_{02}=-\frac1{2K}$ and $C_{10}=C_{21}=\frac1{2K}$. The response peaks near the
endogenous frequency $\sqrt3/(4K)$. At $K=0.4$ and resonance the $x_0$-amplitude per unit
$\varepsilon$ is $0.94$, $0.78$, $1.20$ for $m=0,1,2$. If the clock is only partly synchronized,
with class order parameters $r_\rho=|\langle e^{i\theta_i}\rangle_{V_\rho}|$, the first-order
forcing is multiplied by $r_\rho$. So a macroscopic response exists at first order iff the clock
is synchronized ($r>0$).
\end{theorem}

\begin{proof}
The linearization at $x^{*}$ is $-I+C$, which is hyperbolic and stable for $K>\frac14$. The forcing
enters through $\sigma'(0)/K=1/(4K)$. The periodic solution and its expansion follow from the
implicit function theorem for periodic forcing of a hyperbolic equilibrium. For a partly
synchronized clock, the class average of $\cos\theta_i$ equals $r_\rho\cos\psi_\rho$ to first
order.
\end{proof}

\begin{remark}[Coherence on large complexes]\label{V3P5:rem:coherent}
On the complexes themselves, the class-density standard deviation of the forced logit dynamics
($\varepsilon=0.3$, clock from the flux sector with $\alpha=0.3$) is:
\begin{center}
\small
\begin{tabular}{llccc}
\toprule
& clock & $Y_{24}$ & $Y_{768}$ & $Y_{6144}$\\
\midrule
$K=0.4$, $\omega=1.083$ & synchronized & $0.215$ & $0.110$ & $0.103$\\
& random offsets & $0.201$ & $0.037$ & $0.013$\\
& none & $0.222$ & $0.035$ & $0.013$\\
\midrule
$K=0.1$, $\omega=1.927$ & synchronized & $0.233$ & $0.156$ & $0.148$\\
& random offsets & $0.239$ & $0.041$ & $0.022$\\
& none & $0.247$ & $0.069$ & $0.028$\\
\bottomrule
\end{tabular}
\end{center}
With a synchronized clock the macroscopic oscillation no longer decays with $n$, both above and
below $K_H$. With random offsets it decays like the unforced dynamics. Line (NA3) checks the
$Y_{168}$ case. What produces coherence is the \emph{common clock}. The cycle is entrained,
i.e.\ imposed by an external rhythm whose phases were synchronized separately. It is not an
endogenous macroscopic cycle of the game. The game contributes the resonance and the chirality
dependence of Theorem~\ref{V3P5:thm:entrain}. The ``iff'' of the proposed programme holds in this
sense: coherent macroscopic response iff the clock order parameter is positive. For
heterogeneous clock frequencies that is the classical Kuramoto synchronization condition, which
concerns the clock and not the game.
\end{remark}

\paragraph{Anchors.} A \emph{committed set} $Z$ never changes strategy; it cooperates always.
A chain leader enforcing cooperation in its own chambers is the case $Z=$ the star of a
vertex.

\begin{proposition}[Frontier principle]\label{V3P5:prop:frontier}
Under zero-noise imitation, a committed set changes only the inequalities at its own boundary.
In the thick building ($q=2$, computed in the level-$14$ cover, whose balls of radius $4$ have
$1,14,98,560,2912$ vertices), we find the following.
\begin{enumerate}
\item[(i)] From an anchored ball of radius $1$ (the leader's star), cooperation spreads to the
next sphere iff $r>\frac73,\frac{63}{19},\frac{63}{55}$ (linear, threshold, quadratic).
\item[(ii)] From an anchored ball of radius $2$, it spreads iff $r>\frac{21}{10},3,\frac{21}{20}$.
But the newly reached ball of radius $3$ is not anchored, and it is stable only for
$r>\frac73,\frac{63}{19},\frac{63}{55}$.
\end{enumerate}
So sustained growth beyond an anchor of radius $\le2$ needs $r>\frac{63}{19}$ in the threshold
game, whatever the anchor. In the thin lattice an anchored wall spreads, and every new front is
again a flat wall, stable for $r>\frac95,\frac{18}{11},\frac35$.
\end{proposition}

\begin{proof}
A defector $j$ adjacent to the cooperating set converts iff some cooperating neighbour earns more.
This gives the thresholds $\min_i D_2/(\Phi_i-\Phi_j)$, maximized over the frontier. An
unanchored cooperator keeps its strategy iff it earns more than every defecting neighbour
(Theorem~\ref{V3P5:thm:cluster}). The values are exact computations over the ball (line (NA4)).
\end{proof}

The mechanism-design reading is negative. A leader's enforced contract protects its own
neighbourhood, but the moving front has the geometry of spheres, and spheres in a thick building
have the star thresholds. The absence of flat interfaces is a property of the geometry, and no
finite anchor removes it. Only an \emph{infinite} anchored set shaped like a wall would help, and
thick buildings have no sets with the small boundary of a wall.

\begin{remark}[Anchors at finite noise]\label{V3P5:rem:anchorsmc}
Threshold game, Fermi rule, $K=0.1$. Cooperator fraction among the non-committed vertices:
\begin{center}
\small
\begin{tabular}{lcccc}
\toprule
& $r=2.0$ & $2.25$ & $2.4$ & $2.5$\\
\midrule
$Y_{168}$, no anchor & 0.00 & 0.00 & 0.27 & 0.57\\
$Y_{168}$, leader star ($15$ committed) & 0.00 & 0.00 & 0.65 & 0.82\\
$Y_{168}$, $15$ random committed & 0.00 & 0.00 & 0.68 & 0.80\\
$T_{12}$, no anchor & 0.00 & 0.04 & 0.72 & 0.89\\
$T_{12}$, committed row ($12$) & 0.01 & 0.64 & 0.87 & 0.95\\
$T_{12}$, $12$ random committed & 0.00 & 0.60 & 0.87 & 0.92\\
\bottomrule
\end{tabular}
\end{center}
About $9\%$ of committed vertices raises cooperation near the transition on both geometries. The
leader's star does no better than random commitment. On the thick complex cooperation is still
absent at $r=2.25$, where the thin lattice with a committed row already has $0.64$. Anchors
shift the transition slightly; they do not pin it low. Line (NA5) checks the $r=2.25$ and $2.5$
columns.
\end{remark}

\section{Verification battery}\label{V3P5:sec:battery}

\begin{table}[ht]
\centering\scriptsize
\renewcommand{\arraystretch}{1.2}
\resizebox{\textwidth}{!}{%
\begin{tabular}{lcl}
\toprule
check & complexes & result\\
\midrule
(R) chamber payoff $=$ reduced additive payoff; $w=q+1$ on edges, $0$ elsewhere; $D_2=\kappa(q+1)$ & $Y_{21},Y_{168},T_9$ & numeric $10^{-9}$\\
(Dm) own switch changes own payoff by $D_2(r/3-1)$ & same & numeric\\
(Z) $r_*=3.5$ both ways; margins of stars, chambers, edges; star of $Y_{168}$ at $7/3$; $5/3\notin\{21/(7+\delta)\}$ & $Y_{21},Y_{84},Y_{168},T_{12}$ & exact\\
(MC) Fermi, $K=0.1$: $Y_{168}$, random $14$-regular graph, $T_{12}$ & --- & Monte Carlo (seeded)\\
(P) $\frac{d}{dt}\log V$ identity; $\mu_H$ and frequency; cycle amplitudes; unique equilibrium & --- & numeric\\
(E) $|F|/|V|$, $D_2$, $r_*$ for $q=2,3,4,5,7,8$ & --- & exact\\
(NL1) multilinear coefficients; pairwise iff $f$ affine; $s_js_k$ coefficient $c_2/3$ & $Y_{21}$ & exact\\
(NL2) exact potential; type-stratified Nash window $[\frac32,3]$ & $Y_{21},Y_{168},T_{12}$ & exact\\
(NL3) best cages by enumeration of $2^{21}$ sets; $\max\Phi$ switches at $r=3$ & $Y_{21}$, random & exact\\
(NL4) star cages $=$ closed form & $Y_{168},T_{12}$ & exact\\
(NL5) threshold game, Fermi, $K=0.1$ & $Y_{21},T_{12},Y_{168}$, random & Monte Carlo (seeded)\\
(NC1)--(NC2) circulation $\pm2/7$; strict improvement $6$-cycle & $Y_{21}$ & exact\\
(NC3) Lyapunov function of potential mean fields & --- & numeric\\
(NC4) Hopf points $\frac14$ and $0.170582$; amplitudes; period $2\pi/\sqrt3$; unique equilibrium & --- & numeric\\
(NC5) flat interfaces $18/(3\varphi_3+\varphi_2-\varphi_1)$ & $T_{12}$ & exact\\
(NC6) class-density oscillation shrinks with $n$ & $Y_{24},Y_{168},T_{12},T_{48}$ & Monte Carlo (seeded)\\
(NA1) gauge clock reaches $m^{*}=-\operatorname{round}(\alpha)$ with order $>0.999$ & $Y_{24},Y_{168}$ & numeric\\
(NA2) linear response $X_m(\omega)$ vs forced ODE ($3\%$); resonance & --- & numeric\\
(NA3) synchronized vs random clock vs none & $Y_{168}$ & Monte Carlo (seeded)\\
(NA4) spreading from anchored balls; stability of the radius-$3$ ball; anchored wall & building (level $14$), $T_{24}$ & exact\\
(NA5) committed star vs random vs none; committed row & $Y_{168},T_{12}$ & Monte Carlo (seeded)\\
\bottomrule
\end{tabular}}
\medskip
\caption{Checks run by \texttt{code/v3p5\_verify.py}.}
\label{V3P5:tab:battery}
\end{table}

\section{Scope, residue, corrections}\label{V3P5:sec:scope}

\begin{remark}[Solved, open]\label{V3P5:rem:scope}
\emph{Solved:} the pairwise reduction; dominance and the mean-field threshold; the exact
zero-noise cluster criterion, with its integrality consequence and the absence of cages on
$Y_{21}$; the trichotomy for cyclic replicator dynamics; the Hopf point with mutation.
\emph{Open:} (a) finite-noise thresholds on families of thick complexes of growing size; (b) a
pair or cluster approximation adapted to the $\Atwo$ $1$-skeleton; (c) a proof that the
interior equilibrium with mutation is unique, and the sign of the first Lyapunov coefficient;
(d) genuinely higher-order (nonlinear) chamber games, treated in
Section~\ref{V3P5:sec:nonlinear}. What remains open there is the star formula for $q\ge3$, the
optimal cages on large thick complexes, and finite-noise thresholds along growing families;
(e) three-population versions with a three-dimensional mean field, where tori are possible.
\end{remark}

\begin{remark}[Corrections to the proposed programme]\label{V3P5:rem:corrections}
(1) The chamber public-goods game is a pairwise game on the $1$-skeleton
(Theorem~\ref{V3P5:thm:reduction}). The $\PG(2,q)$ flags enter only through the constant
$q+1$, and there is no separate ``shielding cage''. (2) $r_c=1+2/(q+1)$ is not derived by any
of the available notions. The mean-field threshold is $3$ (Proposition~\ref{V3P5:prop:mf}),
and zero-noise cluster thresholds are $3\kappa/(\kappa+\delta)$, which never equals $5/3$ for
$q=2$ (Theorem~\ref{V3P5:thm:cluster}). The best cage found, the star of $Y_{168}$, needs
$r>7/3$, and $Y_{21}$ has none (Proposition~\ref{V3P5:prop:Y21}). (3) There is no topological
catalysis. Thick complexes behave like random regular graphs, and the thin lattice supports
cooperation better (Remark~\ref{V3P5:rem:mc}). ``$r_c>3$ on regular networks'' is not the
correct baseline (Remark~\ref{V3P5:rem:three}). (4) The three-strategy replicator dynamics has
no limit cycle and no torus (Theorem~\ref{V3P5:thm:rps}). Persistent oscillations arise as a
limit cycle once mutation is added (Theorem~\ref{V3P5:thm:mutation}). (5) $|F|/|V|$ is a
function of $q$ and does not enter the normalized dynamics.
\end{remark}

\section*{Acknowledgement of provenance and caveats}

Unrefereed. The Fermi rule is that of \cite{ST98}. Lattice reciprocity and the $r_c/G$ values
are quoted from \cite{NM92,SPS09}, and the mutation-driven cycle is compared with \cite{Mob10}.
LaSalle's principle \cite{LaSalle60} is classical. The bibliographic data of
\cite{ST98,NM92,SPS09,Mob10} have been checked against Crossref and zbMATH.



